% ==============================================================================
% Section 12: Institutional Capital Adequacy & Model Risk Governance
% Target: Page 14
% ==============================================================================

\section{Institutional Capital Adequacy and Model Risk Governance}
\label{sec:institutional_capital}

\subsection{Supervisory Capital Charge Determination}

Under the Basel Committee's market risk framework (mandated by the Reserve Bank of India for scheduled commercial banks and primary dealers), the minimum regulatory capital requirement for a trading desk operating under the Internal Models Approach is defined as:
\begin{equation}
    C_t = \max \left( \VaR_{t-1}(10\text{d}, 99\%), \; k \cdot \frac{1}{60} \sum_{i=1}^{60} \VaR_{t-i}(10\text{d}, 99\%) \right) + \text{sVaR}_t
    \label{eq:basel_capital_charge}
\end{equation}
where $k = 3 + s$ represents the supervisory market risk multiplier, with $s \in [0, 1]$ being the penalty factor dictated by the Basel Traffic Light backtesting zone:
\begin{equation}
    s = \begin{cases} 0.00, & \text{if } x \le 4 \quad \text{(Green Zone)} \\ 0.40 \text{ to } 0.85, & \text{if } 5 \le x \le 9 \quad \text{(Yellow Zone)} \\ 1.00, & \text{if } x \ge 10 \quad \text{(Red Zone / Revocation)} \end{cases}
    \label{eq:penalty_factor}
\end{equation}

Because our empirical backtesting verified that both the Parametric model ($x = 4$) and Historical Simulation model ($x = 3$) reside safely within the \textbf{Green Zone}, the penalty surcharge is zero ($s = 0$), establishing the baseline multiplier $k = 3.0$.

Applying this regulatory multiplier to our empirical 10-day $99\%$ Historical $\VaR$ (\inr 61,487.89), the minimum baseline regulatory capital charge for holding a single market lot of ITC Limited is:
\begin{equation}
    \text{Regulatory Capital Charge} = 3.0 \times \text{\inr}61,487.89 = \mathbf{\text{\inr}1,84,463.67}
    \label{eq:capital_computed}
\end{equation}
This capital charge represents \textbf{$37.52\%$ of the total market value of the equity position} ($V_0 = \text{\inr}4,91,625.00$). This substantial capital allocation requirement illustrates how regulatory conservatism ensures institutional resilience against unhedged equity holdings.

\subsection{Model Risk Management and Governance Policies}

To mitigate \textbf{model risk}—the potential for financial loss resulting from decisions based on incorrect or misused model outputs—institutional risk committees enforce four mandatory operational controls:

\begin{enumerate}[leftmargin=1.8em, label=(\roman*)]
    \item \textbf{Dynamic Volatility Updating:} Rather than relying on unweighted sample standard deviations, trading desks should deploy exponentially weighted moving average (EWMA, $\lambda = 0.94$) or GARCH(1,1) architectures to capture volatility clustering and reduce exception persistence.
    \item \textbf{Liquidity Horizon Adjustments:} Under FRTB rules, single-stock equity blocks exceeding daily trading participation limits ($>5\%$ of average daily volume) must be scaled to a 20-day or 40-day liquidity horizon.
    \item \textbf{Independent Model Validation:} Risk algorithms must undergo annual independent validation by a quantitative audit team completely segregated from front-office trading desks.
    \item \textbf{Integrated Hedging Requirements:} When capital charges approach institutional thresholds, trading mandates dictate the execution of formal derivative hedging programs (such as Protective Puts or Collars) to truncate tail risk and achieve capital relief.
\end{enumerate}

\clearpage
